\documentclass[10pt]{article}
\usepackage[a4paper,top=14mm,bottom=16mm,left=14mm,right=14mm]{geometry}
\usepackage{fontspec}
\providecommand{\euro}{€}
\usepackage[french]{babel}
\usepackage{librebaskerville}
\usepackage{josefin}
\usepackage{xcolor}
\usepackage{graphicx}
\usepackage{tikz}
\usepackage{multicol}
\usepackage{fontawesome5}
\usepackage{needspace}
\usepackage{fancyhdr}
\usepackage{refcount}
\usepackage[hidelinks]{hyperref}
\makeatletter\newcommand{\pagesuniques}[1]{\def\pu@prec{}\def\pu@sep{}\@for\pu@l:=#1\do{\edef\pu@p{\getpagerefnumber{\pu@l}}\ifx\pu@p\pu@prec\else\pu@sep\pu@p\def\pu@sep{, }\let\pu@prec\pu@p\fi}}\makeatother
\definecolor{titre}{HTML}{F7F1E3}
\definecolor{accent}{HTML}{B5382B}
\definecolor{texte}{HTML}{2B2622}
\definecolor{doux}{HTML}{6A5D50}
\definecolor{carte}{HTML}{E4D3AE}
\definecolor{carteTrait}{HTML}{2F3A36}
\definecolor{cartePoint}{HTML}{B5382B}
\definecolor{fond}{HTML}{F1E6CF}
\definecolor{ardoise}{HTML}{2F3A36}
\colorlet{rubrique}{ardoise}
\pagecolor{fond}
\color{texte}
\setlength{\parindent}{0pt}
\setlength{\parskip}{0pt}
\pagestyle{fancy}\fancyhf{}\renewcommand{\headrulewidth}{0pt}
\fancyfoot[C]{\footnotesize\color{doux}Carte des vins \quad\textperiodcentered\quad \thepage}
\newcommand{\Section}[1]{\clearpage\noindent\colorbox{ardoise}{\parbox{\dimexpr\linewidth-2\fboxsep}{\centering\vspace{2mm}{\sffamily\fontsize{26}{30}\selectfont\bfseries\color{titre}\MakeUppercase{#1}}\vspace{2mm}}}\par\vspace{6mm}}
\newcommand{\Pays}[1]{\Needspace{6\baselineskip}\vspace{4mm}{\sffamily\Large\bfseries\color{ardoise}\MakeUppercase{#1}}\par\vspace{1mm}}
\newcommand{\Region}[1]{\Needspace{5\baselineskip}\vspace{2.5mm}{\sffamily\large\color{accent}#1}\par\vspace{1mm}}
\newcommand{\Appellation}[1]{\Needspace{4\baselineskip}\vspace{1.8mm}{\sffamily\bfseries\color{ardoise}#1}\par\vspace{0.5mm}}
\newcommand{\Bin}[1]{\makebox[2.6em][l]{\footnotesize\color{doux}#1}}
\newcommand{\Vin}[7]{\Needspace{2\baselineskip}\noindent\Bin{#1}#7{\producteur{#2}}\ifx&#3&\else, {\itshape #3}\fi\ {\color{doux}\dotfill}\ #4\ifx&#5&\else\hspace{1.2em}\makebox[3.6em][r]{#5}\fi\par\ifx&#6&\else{\leftskip=2.6em\footnotesize\itshape\color{doux}#6\par}\fi\vspace{0.9mm}}
\newcommand{\Verticale}[4]{\Needspace{3\baselineskip}\noindent\hspace*{2.6em}#4{\producteur{#1}}\ifx&#2&\else, {\itshape #2}\fi\ {\footnotesize\color{doux}(verticale)}\par\ifx&#3&\else{\leftskip=2.6em\footnotesize\itshape\color{doux}#3\par}\fi\vspace{0.3mm}}
\newcommand{\Millesime}[3]{\noindent\Bin{#1}\hspace*{1.2em}#2\ifx&#3&\par\else\ {\color{doux}\dotfill}\ \makebox[3.6em][r]{#3}\par\fi\vspace{0.4mm}}
\newcommand{\coeur}{{\color{accent}\faHeart}\ }
\newcommand{\derniere}{{\footnotesize\color{accent}\faWineBottle}\ }
\newcommand{\producteur}[1]{{\bfseries #1}}
\begin{document}
{\centering{\fontsize{26}{30}\selectfont\color{rubrique} Carte des vins}\\[2mm]{\color{doux} La maison des vendanges --- 26 octobre 2026}\par}\vspace{4mm}
\renewcommand{\Section}[1]{\Needspace{6\baselineskip}\vspace{4mm}{\centering\Large\color{rubrique}#1\par}\vspace{2mm}}
\begin{multicols}{2}
\Section{Au verre et à la carafe}
\Vin{101}{Maison Aubertin \& Fils}{Brut nature, Champagne}{sans année}{18~\euro{}}{Avec : Huîtres, coquillages crus, sushis, poisson cru}{}
\Vin{209}{Domaine du Coteau Fleuri}{Les Silex, Sancerre}{2013}{13~\euro{}}{Avec : Huîtres, coquillages crus, poisson grillé, poisson blanc}{}
\Section{Champagnes et effervescents}
\Pays{France}
\Region{Champagne}
\Appellation{Champagne}
\Vin{101}{Maison Aubertin \& Fils}{Brut nature}{sans année}{90~\euro{}}{Avec : Huîtres, coquillages crus, sushis, poisson cru}{\coeur}\label{l56bouteille}
\Verticale{Maison Aubertin \& Fils}{Brut réserve}{Avec : Homard, langouste, crustacés, foie gras}{\coeur}
\Millesime{102}{2010}{100~\euro{}}\label{l23bouteille}
\Millesime{103}{2019}{85~\euro{}}\label{l24bouteille}
\Verticale{Maison Aubertin \& Fils}{Rosé de saignée}{Avec : Homard, langouste, crustacés, foie gras}{}
\Millesime{104}{2008}{180~\euro{}}\label{l27bouteille}
\Millesime{105}{2024 \derniere}{135~\euro{}}\label{l28bouteille}
\Appellation{Champagne grand cru}
\Verticale{Maison Aubertin \& Fils}{Blanc de blancs}{Avec : Homard, langouste, crustacés, foie gras}{\coeur}
\Millesime{106}{2010 \derniere}{155~\euro{}}\label{l25bouteille}
\Millesime{107}{2022}{165~\euro{}}\label{l26bouteille}
\Section{Vins blancs}
\Pays{France}
\Region{Alsace}
\Appellation{Alsace grand cru}
\Vin{201}{Domaine Weinberg-Hoff}{Grand cru Schlossberg}{2006}{80~\euro{}}{Avec : Huîtres, coquillages crus, poisson grillé, poisson blanc}{}\label{l38bouteille}
\Region{Bordeaux}
\Appellation{Pessac-Léognan}
\Verticale{Clos du Vieux Puits}{}{Avec : Huîtres, coquillages crus, poisson grillé, poisson blanc}{}
\Millesime{202}{2011}{80~\euro{}}\label{l8bouteille}
\Millesime{203}{2012 \derniere}{80~\euro{}}\label{l9bouteille}
\Millesime{204}{2022}{85~\euro{}}\label{l10bouteille}
\Region{Bourgogne}
\Appellation{Chablis premier cru}
\Vin{205}{Domaine Petit-Roc}{Vieilles vignes}{2010}{65~\euro{}}{Avec : Huîtres, coquillages crus, poisson grillé, poisson blanc}{}\label{l21bouteille}
\Appellation{Meursault}
\Vin{206}{Domaine de la Combe Haute}{}{2022}{115~\euro{}}{Avec : Poisson en sauce, au beurre blanc, à la crème, homard, langouste, crustacés}{\coeur}\label{l19bouteille}
\Appellation{Puligny-Montrachet premier cru}
\Vin{207}{Domaine de la Combe Haute}{Les Perrières}{2009}{265~\euro{}}{Avec : Poisson en sauce, au beurre blanc, à la crème, homard, langouste, crustacés}{\derniere}\label{l20bouteille}
\Region{Jura}
\Appellation{Château-Chalon}
\Vin{208}{Domaine des Lauzes}{}{2009}{135~\euro{}}{Avec : Volaille à la crème, aux morilles, comté, beaufort, gruyère (pâtes pressées cuites)}{}\label{l42bouteille}
\Region{Val de Loire}
\Appellation{Sancerre}
\Vin{209}{Domaine du Coteau Fleuri}{Les Silex}{2013}{65~\euro{}}{Avec : Huîtres, coquillages crus, poisson grillé, poisson blanc}{}\label{l36bouteille}
\Region{Vallée du Rhône}
\Appellation{Condrieu}
\Vin{210}{Domaine de la Serre Noire}{}{2006}{115~\euro{}}{Avec : Poisson en sauce, au beurre blanc, à la crème, homard, langouste, crustacés}{\coeur}\label{l32bouteille}
\Pays{Allemagne}
\Region{Moselle}
\Appellation{Mosel}
\Verticale{Weingut Sonnenhang}{Kabinett}{Avec : Huîtres, coquillages crus, poisson grillé, poisson blanc}{}
\Millesime{211}{2016}{50~\euro{}}\label{l51bouteille}
\Millesime{212}{2024}{60~\euro{}}\label{l52bouteille}
\Section{Vins rosés}
\Pays{France}
\Region{Provence}
\Appellation{Bandol}
\Verticale{Domaine du Plateau de Lune}{}{Avec : Poisson grillé, poisson blanc, charcuterie, terrines}{}
\Millesime{401}{2010}{50~\euro{}}\label{l39bouteille}
\Millesime{402}{2011}{60~\euro{}}\label{l40bouteille}
\Millesime{403}{2019}{50~\euro{}}\label{l41bouteille}
\Section{Vins rouges}
\Pays{France}
\Region{Beaujolais}
\Appellation{Morgon}
\Vin{501}{Domaine du Grand Mûrier}{}{2024}{45~\euro{}}{Avec : Charcuterie, terrines, volaille rôtie}{}\label{l22bouteille}
\Region{Bordeaux}
\Appellation{Pauillac}
\Verticale{Château des Trois Chênes}{Grand vin}{Avec : Viande rouge grillée, côte de b\oe{}uf, agneau}{}
\Millesime{502}{2016}{225~\euro{}}\label{l1bouteille}
\Millesime{503}{2019}{195~\euro{}}\label{l3bouteille}
\Verticale{Château des Trois Chênes}{Second vin}{Avec : Viande rouge grillée, côte de b\oe{}uf, agneau}{}
\Millesime{504}{2005}{90~\euro{}}\label{l4bouteille}
\Millesime{505}{2010}{95~\euro{}}\label{l5bouteille}
\Millesime{506}{2020}{90~\euro{}}\label{l6bouteille}
\Appellation{Saint-Émilion grand cru}
\Vin{507}{Château la Garenne}{}{2013}{130~\euro{}}{Avec : Viande rouge grillée, côte de b\oe{}uf, agneau}{}\label{l7bouteille}
\Region{Bourgogne}
\Appellation{Gevrey-Chambertin}
\Verticale{Domaine Marchand-Lavaux}{Les Combes}{Avec : Volaille rôtie, plats mijotés (b\oe{}uf bourguignon, daube, civet)}{}
\Millesime{508}{2013}{160~\euro{}}\label{l14bouteille}
\Millesime{509}{2020 \derniere}{140~\euro{}}\label{l15bouteille}
\Appellation{Vosne-Romanée premier cru}
\Verticale{Domaine Marchand-Lavaux}{Clos de la Fontaine}{Avec : Volaille rôtie, plats mijotés (b\oe{}uf bourguignon, daube, civet)}{}
\Millesime{510}{2005}{335~\euro{}}\label{l16bouteille}
\Millesime{511}{2006}{370~\euro{}}\label{l17bouteille}
\Millesime{512}{2018 \derniere}{275~\euro{}}\label{l18bouteille}
\Region{Languedoc}
\Appellation{Pic-Saint-Loup}
\Verticale{Mas de l'Olivette}{}{Avec : Agneau, plats mijotés (b\oe{}uf bourguignon, daube, civet)}{}
\Millesime{513}{2008 \derniere}{45~\euro{}}\label{l43bouteille}
\Millesime{514}{2012}{45~\euro{}}\label{l44bouteille}
\Millesime{515}{2015}{50~\euro{}}\label{l45bouteille}
\Region{Val de Loire}
\Appellation{Chinon}
\Vin{516}{Domaine des Rochettes}{}{2010}{35~\euro{}}{Avec : Viande rouge grillée, côte de b\oe{}uf, barbecue}{}\label{l37bouteille}
\Region{Vallée du Rhône}
\Appellation{Châteauneuf-du-Pape}
\Vin{517}{Domaine des Galets Roulés}{}{2020}{100~\euro{}}{Avec : Agneau, plats mijotés (b\oe{}uf bourguignon, daube, civet)}{\coeur}\label{l29bouteille}
\Appellation{Côte-Rôtie}
\Vin{518}{Domaine de la Serre Noire}{La Landonne du Bas}{2018}{190~\euro{}}{Avec : Viande rouge grillée, côte de b\oe{}uf, agneau}{\derniere}\label{l31bouteille}
\Pays{Espagne}
\Region{Rioja}
\Appellation{Rioja}
\Vin{519}{Bodegas Peña Alta}{Reserva}{2008}{60~\euro{}}{Avec : Agneau}{}\label{l49bouteille}
\Pays{Italie}
\Region{Piémont}
\Appellation{Barolo}
\Vin{520}{Tenuta Collealto}{}{2017}{165~\euro{}}{Avec : Gibier, champignons, truffe, risotto}{\coeur\derniere}\label{l46bouteille}
\Pays{États-Unis}
\Region{Californie}
\Appellation{Napa Valley}
\Verticale{Ridgeback Cellars}{Estate}{Avec : Viande rouge grillée, côte de b\oe{}uf, barbecue}{}
\Millesime{521}{2013}{210~\euro{}}\label{l53bouteille}
\Millesime{522}{2016 \derniere}{160~\euro{}}\label{l54bouteille}
\Millesime{523}{2022 \derniere}{205~\euro{}}\label{l55bouteille}
\Section{Vins doux et liquoreux}
\Pays{France}
\Region{Bordeaux}
\Appellation{Sauternes}
\Verticale{Château Bel-Orme du Sud}{}{Avec : Foie gras, roquefort, bleus, stilton}{}
\Millesime{601}{2013}{140~\euro{}}\label{l11bouteille}
\Millesime{602}{2017}{110~\euro{}}\label{l12bouteille}
\Region{Val de Loire}
\Appellation{Vouvray}
\Verticale{Domaine du Coteau Fleuri}{}{Avec : Foie gras, cuisine épicée (thaïe, sichuan, indienne)}{}
\Millesime{603}{2008}{70~\euro{}}\label{l33bouteille}
\Millesime{604}{2017}{65~\euro{}}\label{l34bouteille}
\Millesime{605}{2021}{70~\euro{}}\label{l35bouteille}
\Section{Vins mutés}
\Pays{Portugal}
\Region{Douro}
\Appellation{Porto}
\Vin{701}{Quinta do Vale Escuro}{Vintage}{2010}{170~\euro{}}{Avec : Roquefort, bleus, stilton, chocolat}{}\label{l50bouteille}
\Section{Demi-bouteilles}
\Pays{Vins doux et liquoreux}
\Region{Bordeaux}
\Appellation{Sauternes}
\Vin{801}{Château Bel-Orme du Sud}{}{2013}{}{Avec : Foie gras, roquefort, bleus, stilton}{}\label{l11demi}
\Pays{Vins mutés}
\Region{Roussillon}
\Appellation{Banyuls}
\Vin{802}{Mas des Vents Marins}{Rimage}{2021}{}{Avec : Roquefort, bleus, stilton, chocolat}{}\label{l57cinquante}
\Section{Grands formats}
\Pays{Vins rouges}
\Region{Bordeaux}
\Appellation{Pauillac}
\Vin{901}{Château des Trois Chênes}{Grand vin, Magnum}{2017}{}{Avec : Viande rouge grillée, côte de b\oe{}uf, agneau}{}\label{l2magnum}
\end{multicols}
\vfill{\footnotesize\color{doux}\noindent\coeur coup de c\oe{}ur\quad\derniere dernière bouteille\quad BIN : numéro de la bouteille à donner au service.\par}
\end{document}
